\chapter{Farewells}
\label{ch:farewells}
\hlchaptermodality{4}

Telugu, like every Dravidian language, does not say a plain ``goodbye.'' It
promises to return:

\begin{center}
\te{వెళ్ళి వస్తాను.} \quad(\emph{ve\d{l}\d{l}i vast\=anu} --- ``I'll go and come back.'')
\end{center}

% canonical-insertion: TE-S123-letter-dda
\section[ḍa]{\te{డ} — one character, met inside words you already say}

\label{lesson:TE-S123-letter-dda}



\begin{quote}
Before the new one: \te{౫} — what amount does it stand for? And one from further back: \te{ఉ}?

One character this time — and you have been saying it for pages without knowing which mark on the page it was.
\end{quote}

\subsection*{Script you'll notice: \te{డ}}
\textbf{\te{డ}} — \emph{ḍa}.

It is a \textbf{consonant}, and in this script a consonant is never bare: it comes with an \emph{a} already in it. So it is not \emph{ḍ}, it is \textbf{ḍa}.

You met it on an earlier page, in the middle of a word you already say:

\begin{itemize}
\raggedright
  \item \textbf{\te{మీరు} \te{ఎలా} \te{ఉన్నారు}?} — how are you? (respectful)
\end{itemize}

\subsection*{Writing: \te{డ} — copy what you see}
Put your pen on \te{డ} and follow its line. Copy the shape you can see — slowly, and larger than it is printed.

\begin{quote}
This book does not yet tell you \textbf{where to start the character or which way to travel}. That is a real thing, taught with real variation from school to school, and it is not written down here until it can be written down with a source. Copying what is in front of you needs no such source, and it is how the shape gets into your hand in the meantime.
\end{quote}

\subsection*{Your turn}
\begin{itemize}
\raggedright
  \item \emph{Look:} at these words, and find \te{డ} in the ones that have it
\end{itemize}

\begin{quote}
\te{మీరు} \te{ఎలా} \te{ఉన్నారు}?  ·  \te{నమస్కారం}
\end{quote}

\begin{itemize}
\raggedright
  \item \emph{Trace it:} \te{డ} three times, saying \emph{ḍa} as you finish each one
  \item \emph{Look:} back at any page of this chapter and find \te{డ} once more
\end{itemize}

\begin{tcolorbox}[breakable,colback=teal!4,colframe=teal!35!black,title={Before you move on}]
Which character is this — \te{డ}? What sound does it carry? (\emph{\textbf{ḍa}}.) Name one word you already say that contains it.

One more, from much earlier: \te{న} — what does it do?
\end{tcolorbox}

\section{\te{వెళ్ళు} \emph{(ve\d{l}\d{l}u)} --- ``go,'' and \te{వచ్చు} (come)}
\label{lesson:vellu}

\begin{sounds}
\te{వె} (\emph{ve}) $+$ \te{ళ్ళు} (\emph{\d{l}\d{l}u}, doubled retroflex \te{ళ}
\emph{\d{l}}) $\to$ \te{వెళ్ళు}. Its partner \te{వచ్చు} (\emph{vaccu}, ``come'').
\end{sounds}

\begin{cousinweb}
\te{వెళ్ళు} (\emph{ve\d{l}\d{l}u}, ``to go'') and \te{వచ్చు} (\emph{vaccu}, ``to
come'') are native Dravidian verbs. As commands they are whole words:
\emph{ve\d{l}\d{l}u!} (``go!''), \emph{r\=a!} (``come!''). You need both, because
Telugu's goodbye is ``I go, and I \emph{come back}.''
\end{cousinweb}

% canonical-insertion: TE-S145-digit-6
\section[6]{\te{౬} — the Telugu numeral for 6}

\label{lesson:TE-S145-digit-6}



\begin{quote}
Before the new one: \te{డ} — what does it do? And one from further back: \te{౪}?

One numeral this time, and it is read the way every numeral is read: as an amount.
\end{quote}

\subsection*{Script you'll notice: \te{౬}}
\textbf{\te{౬}} is \textbf{six}.

Six is where the second half begins, and where a warning belongs: \textbf{\te{౬} is not the 6 you write, and it is not an E}. It is its own character, and the only way it becomes readable is by meeting it often enough that the resemblance stops mattering.

\subsection*{Writing: \te{౬} — copy what you see}
Put your pen on \te{౬} and follow its line. Copy the shape you can see — slowly, and larger than it is printed.

\begin{quote}
This book does not yet tell you \textbf{where to start the character or which way to travel}. That is a real thing, taught with real variation from school to school, and it is not written down here until it can be written down with a source. Copying what is in front of you needs no such source, and it is how the shape gets into your hand in the meantime.
\end{quote}

\subsection*{Your turn}
\begin{itemize}
\raggedright
  \item \emph{Look:} at this row and say the value of each numeral in it, left to right
\end{itemize}

\begin{quote}
\te{౬}  ·  \te{౨}  ·  \te{౫}
\end{quote}

\begin{itemize}
\raggedright
  \item \emph{Trace it:} \te{౬} three times, saying \textquotedblleft{}6\textquotedblright{} as you finish each one
  \item \emph{Look:} for \te{౬} on any page of a Telugu book, calendar or doorplate you can find
\end{itemize}

\begin{tcolorbox}[breakable,colback=teal!4,colframe=teal!35!black,title={Before you move on}]
Which numeral is this — \te{౬}? What amount does it stand for? (\te{౬} is 6.)

One more, from much earlier: \te{◌ు} — what does it do?
\end{tcolorbox}

\section{\te{వెళ్ళి వస్తాను} \emph{(ve\d{l}\d{l}i vast\=anu)} --- ``I'll go and come back''}
\label{lesson:velli-vastaanu}

\te{వెళ్ళి} (\emph{ve\d{l}\d{l}i}, ``having gone'') $+$ \te{వస్తాను}
(\emph{vast\=anu}, ``I come'') --- literally ``\textbf{having gone, I come
[back]},'' i.e. ``I'll go and come back.'' A Telugu speaker never signs off with a
bare ``I'm leaving.'' The warm reply is \te{వెళ్ళి రండి} (\emph{ve\d{l}\d{l}i
ra\d{n}\d{d}i}, ``go and come [back]'').

\subsection*{Across the family --- the same idea, five ways}
\begin{center}
\begin{tabular}{@{}ll@{}}
\toprule
Language & ``go and come back'' (goodbye) \\
\midrule
\textbf{Telugu} & \emph{ve\d{l}\d{l}i vast\=anu} \\
Tamil & \emph{p\=oy varugi\d{r}\=e\d{n}} \\
Kannada & \emph{h\=ogi barutt\=ene} \\
Malayalam & \emph{p\=oyi var\=a\d{m}} \\
\midrule
Bengali (Indo-Aryan) & \emph{\=ashi} (``I come'') \\
\bottomrule
\end{tabular}
\end{center}
The whole south, and beyond, prefers a promise of return to a blunt farewell.

\section{\te{రేపు కలుద్దాం} \emph{(r\=epu kalud\d{d}\=am)} --- ``see you tomorrow''}
\label{lesson:reepu-kaluddaam}

\begin{sounds}
\te{రే} (\emph{r\=e}) $+$ \te{పు} (\emph{pu}) $\to$ \te{రేపు}. \te{కలుద్దాం}
(\emph{kalud\d{d}\=am}) $=$ \emph{kalu} (``meet'') $+$ \emph{-d\d{d}\=am}
(``let's'').
\end{sounds}

\begin{cousinweb}
\te{రేపు} (\emph{r\=epu}, ``tomorrow'') is native; \te{కలు} (\emph{kalu}, ``to
meet, join'') gives \emph{kalud\d{d}\=am} (``let's meet'') --- ``tomorrow, let's
meet.'' The \te{-దాం} (\emph{-d\d{d}\=am}) ending is Telugu's ``let's \ldots'':
\emph{ve\d{l}\d{l}d\d{d}\=am} (``let's go''), \emph{tind\=am} (``let's eat'').
\end{cousinweb}

% canonical-insertion: TE-S113-letter-pa
\section[pa]{\te{ప} — one character, met inside words you already say}

\label{lesson:TE-S113-letter-pa}



\begin{quote}
Before the new one: \te{౬} — what amount does it stand for? And one from further back: \te{ర}?

One character this time — and you have been saying it for pages without knowing which mark on the page it was.
\end{quote}

\subsection*{Script you'll notice: \te{ప}}
\textbf{\te{ప}} — \emph{pa}.

It is a \textbf{consonant}, and in this script a consonant is never bare: it comes with an \emph{a} already in it. So it is not \emph{p}, it is \textbf{pa}.

You already say these, and every one of them has \te{ప} somewhere inside it:

\begin{itemize}
\raggedright
  \item \textbf{\te{పేరు}} \emph{pēru} — name
  \item \textbf{\te{పరవాలేదు}} \emph{paravālēdu} — it's okay / no problem / you're welcome
  \item \textbf{\te{రేపు} \te{కలుద్దాం}} \emph{rēpu kaluddām} — see you tomorrow
\end{itemize}

\subsection*{Writing: \te{ప} — copy what you see}
Put your pen on \te{ప} and follow its line. Copy the shape you can see — slowly, and larger than it is printed.

\begin{quote}
This book does not yet tell you \textbf{where to start the character or which way to travel}. That is a real thing, taught with real variation from school to school, and it is not written down here until it can be written down with a source. Copying what is in front of you needs no such source, and it is how the shape gets into your hand in the meantime.
\end{quote}

\subsection*{Your turn}
\begin{itemize}
\raggedright
  \item \emph{Look:} at these words, and find \te{ప} in the ones that have it
\end{itemize}

\begin{quote}
\te{పేరు}  ·  \te{పరవాలేదు}  ·  \te{నమస్కారం}
\end{quote}

\begin{itemize}
\raggedright
  \item \emph{Trace it:} \te{ప} three times, saying \emph{pa} as you finish each one
  \item \emph{Look:} back at any page of this chapter and find \te{ప} once more
\end{itemize}

\begin{tcolorbox}[breakable,colback=teal!4,colframe=teal!35!black,title={Before you move on}]
Which character is this — \te{ప}? What sound does it carry? (\emph{\textbf{pa}}.) Name one word you already say that contains it.

One more, from much earlier: \te{౧} — what amount does it stand for?
\end{tcolorbox}

\section[Malli kaluddam]{\te{మళ్ళీ కలుద్దాం} \emph{(ma\d{l}\d{l}\=\i\ kalud\d{d}\=am)} --- ``we'll meet again''}
\label{lesson:malli-kaluddaam}

\begin{cousinweb}
\te{మళ్ళీ కలుద్దాం} $=$ \te{మళ్ళీ} (\emph{ma\d{l}\d{l}\=\i}, ``again'') $+$
\emph{kalud\d{d}\=am} (``let's meet''). Both are native Dravidian: \emph{ma\d{l}\d{l}\=\i}
(from \emph{maral}, ``to return'') and \emph{kalu} (``meet''). Where Tamil's
``we'll meet again'' reaches for the Sanskrit loan \emph{sandi} to say ``meet,''
Telugu keeps the native \emph{kalu} --- the same phrase, from different stock.
\end{cousinweb}

% canonical-insertion: TE-S146-digit-7
\section[7]{\te{౭} — the Telugu numeral for 7}

\label{lesson:TE-S146-digit-7}



\begin{quote}
Before the new one: \te{ప} — what does it do? And one from further back: \te{౫}?

One numeral this time, and it is read the way every numeral is read: as an amount.
\end{quote}

\subsection*{Script you'll notice: \te{౭}}
\textbf{\te{౭}} is \textbf{seven}.

\textbf{\te{౬}} and \textbf{\te{౭}} sit next to each other in value and are easy to swap on a first reading. Look at the pair below until you can say which is which without counting up to it.

\subsection*{Writing: \te{౭} — copy what you see}
Put your pen on \te{౭} and follow its line. Copy the shape you can see — slowly, and larger than it is printed.

\begin{quote}
This book does not yet tell you \textbf{where to start the character or which way to travel}. That is a real thing, taught with real variation from school to school, and it is not written down here until it can be written down with a source. Copying what is in front of you needs no such source, and it is how the shape gets into your hand in the meantime.
\end{quote}

\subsection*{Your turn}
\begin{itemize}
\raggedright
  \item \emph{Look:} at this row and say the value of each numeral in it, left to right
\end{itemize}

\begin{quote}
\te{౭}  ·  \te{౬}  ·  \te{౪}
\end{quote}

\begin{itemize}
\raggedright
  \item \emph{Trace it:} \te{౭} three times, saying \textquotedblleft{}7\textquotedblright{} as you finish each one
  \item \emph{Look:} for \te{౭} on any page of a Telugu book, calendar or doorplate you can find
\end{itemize}

\begin{tcolorbox}[breakable,colback=teal!4,colframe=teal!35!black,title={Before you move on}]
Which numeral is this — \te{౭}? What amount does it stand for? (\te{౭} is 7, \te{౬} is 6.)

One more, from much earlier: \te{◌్} — what does it do?
\end{tcolorbox}

\section{The farewells}
\label{lesson:te-c04-practice}

\begin{center}
\begin{tabular}{@{}ll@{}}
\toprule
Telugu & English \\
\midrule
\te{వెళ్ళి వస్తాను.} & I'll go and come back. (goodbye) \\
\te{వెళ్ళి రండి.} & Go and come back. (the reply) \\
\te{రేపు కలుద్దాం.} & See you tomorrow. \\
\te{మళ్ళీ కలుద్దాం.} & We'll meet again. \\
\bottomrule
\end{tabular}
\end{center}

Roots banked: \emph{ve\d{l}\d{l}u}/\emph{vaccu} (go/come), \emph{r\=epu}
(``tomorrow''), \emph{kalu} (``meet,'' native --- where Tamil borrowed
\emph{sandi}), and \emph{ma\d{l}\d{l}\=\i} (``again''). Next chapter: the first
real verbs --- \emph{n\=enu telugu m\=a\d{t}l\=a\d{d}at\=anu} --- and Telugu's
sentences begin to move.
